% called by main.tex
%
\chapter{Introduction}
\label{ch::chapter1}

Esta plantilla sigue el formato de tesis de la \textit{Universidad Politécnica de Madrid} (UPM), que se encuentra disponible\href{https://www.upm.es/Estudiantes/Estudios_Titulaciones/Estudios_Doctorado/Tesis/NormasRedaccionTesis}{aquí}.


En el siguiente capítulo se proporciona información acerca de las reglas parar preparar el documento de la tesis, así como la explicación del uso de esta plantilla.

\section{Motivación}

Los modelos de lenguaje de gran tamaño han demostrado una capacidad muy elevada para generar texto a partir de instrucciones en lenguaje natural. Sin embargo, esa capacidad resulta mucho menos fiable cuando la salida esperada no es texto libre, sino una estructura formal que debe respetar simultáneamente restricciones sintácticas, jerárquicas y semánticas. Este problema aparece con claridad en la generación de manifiestos YAML de Kubernetes, donde una salida aparentemente razonable desde el punto de vista lingüístico puede ser, sin embargo, inválida, inconsistente o directamente inutilizable. En este contexto, no basta con que el modelo ``entienda'' el prompt: debe además construir una representación estructurada correcta y mantener la coherencia entre las distintas partes del documento. Esta es precisamente la motivación técnica central del presente trabajo, que se enmarca en el estudio de la generación estructurada de YAML a partir de lenguaje natural.

El interés del problema surge, en parte, de las limitaciones que presentan las soluciones generalistas actuales. Un LLM convencional, utilizado únicamente mediante prompting o incluso tras un ajuste supervisado básico, genera la salida de manera autoregresiva token a token. Es decir, su mecanismo natural consiste en producir una secuencia plana de tokens, uno detrás de otro. Este mecanismo funciona bien cuando el objetivo es producir texto natural, pero no garantiza por sí mismo que la secuencia generada respete una estructura de árbol, que las claves aparezcan en el contexto adecuado o que las dependencias internas del documento sean coherentes. En dominios técnicos, pequeños errores locales, como una indentación incorrecta, una clave mal anidada o una referencia a un recurso inexistente, pueden invalidar completamente el resultado. Además, una parte importante de los fallos no proviene solo de una mala comprensión del prompt, sino de la dificultad de serializar correctamente la información en un formato fuertemente restringido, lo que evidencia que el problema no puede tratarse como una simple tarea de generación de texto libre.


Frente a estas limitaciones, este trabajo propone abordar el problema desde una perspectiva más controlada. La idea central es que la calidad de la generación estructurada puede mejorar si, además del ajuste fino supervisado, se incorporan mecanismos explícitos de control de estructura y de validación automática, en lugar de confiar únicamente en la generación libre del modelo. Para ello, se plantea separar conceptualmente dos dimensiones que en muchos enfoques generalistas aparecen mezcladas: por un lado, la extracción del contenido semántico relevante del prompt y, por otro, la organización formal de dicho contenido en una estructura YAML correcta. Sobre esta base, el proyecto no se limita a preguntar si un modelo puede escribir un manifiesto que parezca correcto, sino que intenta estudiar de qué manera puede representarse y controlarse la jerarquía que hay detrás de ese manifiesto.

Viéndolo de esta forma, el problema puede entenderse desde dos ángulos complementarios. El primero consiste en tratar la estructura como una propiedad de la superficie generada: el modelo produce una representación estructurada serializada, un parser determinista reconstruye el YAML final a partir de ella y, en una fase posterior, podrían utilizarse preferencias automáticas o señales de recompensa para favorecer las salidas más válidas. Esta línea conecta con la motivación de los métodos de alineamiento, pero sin asumir que exista necesariamente un pipeline completo de RLHF con preferencias humanas. El segundo ángulo consiste en tratar el YAML como una estructura de árbol proyectada sobre una secuencia de líneas. En ese caso, el orden de generación proporciona de forma natural la posición de cada entrada, pero no proporciona de manera fiable su altura dentro del árbol. Esa altura se representará en este trabajo mediante la variable \texttt{level}, que codifica el nivel jerárquico de cada línea YAML.

La elección de este tema también está motivada por el interés de trabajar sobre un caso de estudio técnicamente relevante y al mismo tiempo evaluable de forma formal. Kubernetes resulta especialmente apropiado porque permite combinar lenguaje natural, estructuras jerárquicas y reglas verificables en un dominio real, ampliamente utilizado y con una semántica técnica suficientemente rica. En este sentido, el proyecto no solo persigue mejorar la capacidad de un modelo para producir archivos YAML, sino también construir una metodología reproducible para estudiar cómo se comportan los LLMs cuando la tarea exige algo más que fluidez textual.

Por último, existe una motivación personal que también resulta importante en la elección del tema. Más allá del interés por los modelos de lenguaje y por las tareas de NLP, este trabajo nace de la intención de abordar un problema asociado habitualmente a sistemas de tipo ``caja negra'' desde una perspectiva más analítica y, en cierta medida, más matemática. En lugar de limitarse a observar el comportamiento externo del modelo, el objetivo es estudiar el problema mediante métricas objetivas, restricciones formales, análisis estructural del dataset y experimentos que permitan aislar con mayor claridad dónde aparecen los errores y por qué aparecen. Esta motivación conecta con la idea de trasladar herramientas de análisis más rigurosas a un ámbito en el que con frecuencia predominan aproximaciones muy empíricas o poco interpretables. De este modo, el proyecto no solo pretende aportar una solución práctica a una tarea de generación estructurada, sino también servir como marco para estudiar con mayor profundidad el comportamiento de los LLMs en un problema concreto, técnico y medible.

\section{Descripción}

Este Trabajo Fin de Máster se centra en el estudio de la generación estructurada de configuraciones técnicas a partir de instrucciones en lenguaje natural mediante modelos de lenguaje de gran tamaño. En particular, el proyecto aborda como caso de estudio la generación automática de manifiestos de Kubernetes, un dominio especialmente adecuado por combinar una estructura jerárquica bien definida con restricciones sintácticas y semánticas que pueden validarse, al menos parcialmente, de forma automática.

A diferencia de las tareas de generación de texto libre, en este problema no basta con que el modelo produzca una salida lingüísticamente plausible. La salida debe respetar simultáneamente varios niveles de corrección: debe ser sintácticamente válida como YAML, mantener una estructura jerárquica coherente, utilizar claves y valores compatibles con el dominio y reflejar de forma fiel los requisitos expresados en el prompt. Además, en el caso concreto de Kubernetes, la calidad de la generación depende también de relaciones semánticas entre distintos componentes, como la correspondencia entre los selectores de un \texttt{Service} y las etiquetas de los pods que debe exponer, la presencia de campos básicos como \texttt{apiVersion}, \texttt{kind} o \texttt{metadata.name}, la definición correcta de contenedores e imágenes, o la coherencia entre volúmenes declarados y puntos de montaje usados por los contenedores.

El objetivo del proyecto es diseñar y evaluar un sistema capaz de transformar descripciones en lenguaje natural en manifiestos YAML estructurados, válidos y semánticamente consistentes dentro del dominio de Kubernetes. Para ello, el trabajo parte de la idea de que el problema debe abordarse no como una mera tarea de generación autoregresiva token a token, sino como una tarea en la que el modelo debe organizar el contenido del prompt dentro de una estructura formal. Viéndolo de esta forma, podemos entender este tipo de ficheros YAML como estructuras de árbol, donde cada entrada indentada ocupa una posición concreta dentro del documento final. Esta separación resulta especialmente relevante porque muchos de los errores observables en este tipo de tareas no provienen únicamente de una mala interpretación del prompt, sino también de la dificultad del modelo para serializar correctamente la información en un formato fuertemente restringido.

Metodológicamente, el proyecto se apoyará en un modelo base autoregresivo de tipo \textit{decoder-only}, adaptado mediante técnicas de ajuste fino eficiente, como LoRA, sobre un conjunto de pares \textit{prompt--salida}. A partir de esa base, se compararán dos estrategias principales para estudiar la robustez estructural del sistema. La primera, denominada en el repositorio \texttt{serialized\_sft}, mantendrá la jerarquía en la superficie textual de salida: el modelo generará una representación serializada en la que cada línea contiene tanto el contenido YAML como su nivel jerárquico, y un parser determinista se encargará de reconstruir el manifiesto final sin inventar claves ni reparar errores semánticos de forma oculta.

La segunda estrategia, denominada \texttt{two\_head\_sft}, parte de una intuición distinta. Si el YAML se interpreta como un árbol proyectado sobre una secuencia de líneas, el orden de los nodos queda dado por el propio orden de generación del modelo. Lo que no queda determinado de forma fiable es la profundidad de cada entrada. Por ello, esta rama propone separar la predicción del contenido de la predicción de la jerarquía: el modelo generará el texto de cada línea, mientras que un cabezal estructural explícito predecirá el valor \texttt{level}. La comparación entre ambas ramas permitirá estudiar si la jerarquía debe aprenderse como una parte más de la secuencia generada o si conviene modelarla como una señal estructural diferenciada.


La evaluación del sistema se diseñará específicamente para el problema de generación estructurada. Por ello, no se tomarán como referencia principal métricas textuales superficiales, ya que dos configuraciones pueden diferir léxicamente y, sin embargo, ser funcionalmente equivalentes, o parecer similares en texto y ser estructuralmente incorrectas. En su lugar, se definirán métricas en varios niveles: métricas sintácticas, como el porcentaje de salidas YAML parseables; métricas estructurales, como la reconstrucción segura desde bloques, la coincidencia de niveles \texttt{level} o la consistencia de la jerarquía; métricas de dominio Kubernetes, como la presencia de campos requeridos y la coherencia entre selectores, etiquetas, contenedores, puertos o volúmenes; y métricas de adecuación al prompt, destinadas a medir si la configuración generada refleja correctamente la intención del usuario. Esta evaluación se complementará con un análisis sistemático de errores para identificar patrones recurrentes y comprender mejor las limitaciones de cada estrategia.

\section{Objetivos}

El objetivo general de este Trabajo Fin de Máster es diseñar y evaluar un sistema basado en modelos de lenguaje capaz de generar manifiestos YAML de Kubernetes a partir de instrucciones en lenguaje natural, garantizando no solo la validez sintáctica de la salida, sino también su coherencia estructural y semántica respecto al prompt de entrada.

A partir de este objetivo general, se plantean los siguientes objetivos específicos:

\begin{itemize}
    \item Analizar y caracterizar el dataset disponible, estudiando su tamaño, diversidad, complejidad estructural y posibles limitaciones para el entrenamiento del modelo.
    
    \item Diseñar un conjunto de métricas objetivas que permitan evaluar de forma rigurosa la calidad de las salidas generadas, considerando parseabilidad YAML, consistencia estructural, adecuación al prompt y validez de dominio Kubernetes.
    
    \item Definir una representación intermedia basada en bloques de línea, donde cada unidad contenga el texto de la línea y su nivel jerárquico \texttt{level}, de forma que la estructura pueda reconstruirse y evaluarse de manera explícita.
    
    \item Adaptar un modelo base de lenguaje al dominio de generación de configuraciones de Kubernetes mediante técnicas de ajuste fino eficientes.
    
    \item Comparar una estrategia de ajuste supervisado que serializa el nivel jerárquico como texto, \texttt{serialized\_sft}, con una estrategia que incorpora un cabezal estructural explícito para predecir dicho nivel, \texttt{two\_head\_sft}.
    
    \item Evaluar experimentalmente el sistema propuesto sobre un conjunto de prueba representativo, identificando fortalezas, limitaciones y patrones de error tanto en el contenido generado como en la jerarquía predicha.
    
    \item Extraer conclusiones sobre la utilidad de los mecanismos de control estructural en tareas de generación técnica, y sobre la conveniencia de tratar la jerarquía como una parte más de la superficie textual o como una señal estructural diferenciada.
\end{itemize}
